% Ad Familiares 2.7 --- headnote
% ad-familiares-02-07, 19 December 51 BC, from camp at Pindenissus

\textit{Cicero to C. Scribonius Curio, written from the field
before Pindenissus on the fourteenth day before the Kalends of
January, 51 BC (Perseus: \textit{in castris ad Pindenessum a d.
xiv K. Ianuar. a. 703}) --- i.e. 19 December, only a few days
before the siege of the Eleutherocilician stronghold would come
to its end on the day of the Saturnalia. The salutation
\textit{M. Cicero imp. s. d. C. Curioni tr. pl.} marks two
political facts at once: Cicero has been acclaimed imperator by
his army after the Amanus campaign, and Curio has at last
entered the tribunate of the plebs that the previous letters
of book 2 had taken as the still-distant prospect. The
congratulation is therefore late, and Cicero opens with the
defence of a man who hears things on a six weeks' delay from the
edge of the empire. The seriousness of the letter is the
seriousness, by now familiar, of Cicero's whole Curio
correspondence: this is the elder man writing to a brilliant
younger one whom he has long claimed as a pupil, with the
flattery of pupillage and the steadiness of one who would like
to control where he can only advise.}

\textit{The set-piece is the moral-tutor refrain at the heart of
\S 1--2: speak with yourself, take yourself into counsel, listen
to yourself, obey yourself. Cicero presses it twice, once as
opening exhortation and once as recapitulation, and in between
he turns to the broader theme of how Curio has come into office:
not stumbled into the moment but chosen it, brought his
tribunate forward into the very crisis of affairs. He
catalogues, with rhetorical question, what the young tribune
must reckon with --- the force of times, the variety of things,
the uncertainty of outcomes, the pliability of wills, the snares
and emptiness of public life. The flattery is open: Curio does
not need outside counsel, and yet (the immortal gods are
invoked) Cicero wishes he were there as spectator or sharer or
ally. \S 3 sets up the next dispatch and points Curio to a
separate letter sent with his freedman Thraso about his
priesthood. \S 4 is the substantive request --- the standing
refrain of every letter Cicero is writing from Cilicia: do not
let my term be extended. The wording is calibrated to Curio's
new office: he asks now not from one who was merely a senator,
but from a tribune of the plebs, and from Curio the tribune ---
not that something new be decreed (the harder thing) but that
nothing new be decreed, that the senatus consultum and the laws
be defended, and that the condition under which Cicero set out
remain his. The whole letter is the older man's careful work of
investing in the young one before the young one's own course is
fixed; within a year and a half Curio will have crossed to
Caesar and be on the way to the death in Africa that begins the
civil war.}
